\documentclass[12pt]{article}
\usepackage[a4paper, total={6.5in, 9.5in}]{geometry}
\usepackage{graphicx}			
\usepackage{amssymb}
\usepackage{amsmath}			
\usepackage{fancyhdr}
\usepackage{enumerate}       
\usepackage{dirtytalk} 
\usepackage[english]{babel}
\usepackage{amsthm}
\usepackage{xcolor}
\usepackage[shortlabels]{enumitem}
\setlength{\parindent}{0pt}
\setlength{\parskip}{0.8em}

\newtheorem{theorem}{Theorem}[section]
\newtheorem{corollary}{Corollary}[theorem]
\newtheorem{lemma}[theorem]{Lemma}
\newtheorem{definition}[theorem]{Definition}


\begin{document}
% \renewcommand{\thesection}{Exercise \arabic{section}}

\setcounter{section}{4} % Sets the section counter to 4, so the next section is 5
\renewcommand{\thesection}{\arabic{section}}


\parindent = 0pt


      {\bf Analysis 1 HW 2}


%==================

\section{}

 {\bf Definitions}\\
$A$ is a nonempty set of real numbers bounded below. \\
$-A = \{-x \ |\  x \in A\}$. \\

{\bf To prove:}
$\text{inf}(A) = -\text{sup}(-A)$.

\begin{proof}
      Since $A$ is nonempty and bounded below, $y = \text{inf}(A)$ exists. It follows that
      \begin{align*}
                        & y \le x \ \ \forall \, x \in A     \\
            \implies \  & -x \le -y \ \ \forall \, x \in A   \\
            \implies \  & -x \le -y \ \ \forall \, -x \in -A
      \end{align*}
      $\therefore$ $-y$ is an upper bound of $-A$.\\ \\
      Also,
      \begin{align*}
                        & \forall \, x > y \ \ \exists \,x' \in A \text{ such that } x' < x       \\
            \implies \  & \forall \, -x < -y \ \ \exists \,-x' \in -A \text{ such that } -x' > -x
      \end{align*}
      $\therefore$ $-y$ is the smallest upper bound of $-A$.\\
      $-\text{inf}(A) = \text{sup}(-A)$.
\end{proof}

\section{}
 {\bf Definitions}\\
$b \in \mathbb{R}, \  b > 1$.\\
$m, n, p, q \in \mathbb{Z}, \  n>0, \  q>0$.\\
$r \in \mathbb{Q},  \ r = m/n = p/q$.\\
$a^{1/n} \text{ is defined as the positive solution to the equation } y^n = a$ (Rudin, 1.21).




\begin{lemma}
      \label{distributeRoot}
      $(a_1a_2\cdots a_k)^{1/n} = a_1^{1/n}a_2^{1/n}\cdots a_k^{1/n}$
      \begin{proof}
            From the corollary of 1.21, we know that $(a_1a_2)^{1/n} = a_1^{1/n}a_2^{1/n}$.
            Assume the result holds for some $k \ge 2$. Then, $(a_1a_2\cdots a_{k+1})^{1/n} = (a_1a_2\cdots a_{k})^{1/n}a_{k+1}^{1/n}
                  = a_1^{1/n}a_2^{1/n}\cdots a_{k+1}^{1/n}$.
            By the principle of mathematical induction, the lemma is true for all $k \ge 2$.
      \end{proof}
\end{lemma}
\begin{lemma}
      \label{revExpOrder}
      $(b^m)^{1/n} = (b^{1/n})^{m}$
\end{lemma}
\begin{proof}
      If $m\ge 0$ ,
      \begin{align*}
            (b^m)^{1/n} = \Bigg(\prod_{i = 1}^{m}b\Bigg)^{1/n} = \prod_{i = 1}^{m}b^{1/n} = (b^{1/n})^{m}
      \end{align*}
      from \ref{distributeRoot}. A similar
      argument works for $m < 0$.
\end{proof}
\begin {enumerate} [(a)]
\item
      {\bf To prove:}
$(b^m)^{1/n} = (b^p)^{1/q}$.

\begin{proof}
      We know that $mq = pn$. Thus, $b^{mq} = b^{pn}$ or $(b^m)^q = (b^p)^n$. It follows that
      \begin{align*}
                        & ((b^m)^q)^{1/n} = b^p     & \text{Rudin}, 1.21 \\
            \implies \  & ((b^m)^{1/n})^{q} = b^p   & \ref{revExpOrder}  \\
            \implies \  & (b^m)^{1/n} = (b^p)^{1/q} & \text{Rudin}, 1.21
      \end{align*}
\end{proof}
\begin{definition}
      \label{rationalExpDef}
      $b^r\equiv (b^m)^{1/n}$.
\end{definition}
\item
      {\bf To prove:}
$b^{r+s} = b^rb^s, \ \ r,s \in \mathbb{Q}$.

\begin{proof}
      Let $s = c/d, \ \  c, d \in \mathbb{Z},\ \  d>0$. Then,
      \begin{align*}
                 & b^{r+s}                                 &                      \\
            = \  & (b^{md+cn})^{1/{nd}}                    & \ref{rationalExpDef} \\
            = \  & (b^{md}b^{cn})^{1/{nd}}                 &                      \\
            = \  & (b^{md})^{1/{nd}}\cdot(b^{cn})^{1/{nd}} & \ref{distributeRoot} \\
            = \  & (b^{m})^{1/{n}}\cdot(b^{c})^{1/{d}}     & 6(a)                 \\
            = \  & b^rb^s                                  & \ref{rationalExpDef}
      \end{align*}
\end{proof}
\item


\begin{lemma}
      \label{ineqMul}
      If $x, y, z, w \in \mathbb{R}^+,\  x\le y,\ z \le w$, then $xz\le yw$.
\end{lemma}
\begin{proof}
      If $x<y$ and $z<w$,
      \begin{align*}
                           & x<y \implies zx<zy & \text{Rudin, 1.18(b)} \\
                           & z<w \implies yz<yw & \text{Rudin, 1.18(b)} \\
            \therefore\ \  & xz<yw              & \text{Rudin, 1.5(ii)}
      \end{align*}
      If $x=y$ and $z<w$, $xz<yw$ from Rudin, 1.18(b).\\
      If $x<y$ and $z=w$, $xz<yw$ from Rudin, 1.18(b).\\
      If $x=y$ and $z=w$, $xz=yw$.
\end{proof}


\begin{lemma}
      \label{rationalExpComp}
      If $r, s \in \mathbb{Q}, \ r > s$ and $  b \in \mathbb{R}, \  b > 1$
      then $b^r > b^s$.
\end{lemma}

\begin{proof}
      Let $s = c/d, \ \  c, d \in \mathbb{Z},\ \  d>0$. Then, $md > cn$, which implies
      $b^{md} > b^{cn}$. For the sake of contradiction, assume $(b^{md})^{1/nd} \le (b^{cn})^{1/nd}$.
      Then,
      \begin{align*}
                        & \prod_{i = 1}^{nd}(b^{md})^{1/nd} \le \prod_{i = 1}^{nd}(b^{cn})^{1/nd} & \ref{ineqMul}        \\
            \implies \  & ((b^{md})^{nd})^{1/nd} \le ((b^{cn})^{nd})^{1/nd}                       & \ref{distributeRoot} \\
            \implies \  & b^{md} \le b^{cn}                                                       & \text{Rudin, 1.21}
      \end{align*}
      which contradicts $b^{md} > b^{cn}$. Thus,
      \begin{align*}
                        & (b^{md})^{1/nd} > (b^{cn})^{1/nd} &                      \\
            \implies \  & (b^m)^{1/n} > (b^c)^{1/d}         & \text{6(a)}          \\
            \implies \  & b^r > b^s                         & \ref{rationalExpDef}
      \end{align*}
\end{proof}

This lemma can be extended to say that if $r \ge s$ then $b^r \ge b^s$, since $r = s$
implies $b^r = b^s$.

\begin{definition}\label{BxDef}
      $B(x) \equiv \{b^t \ | \  t\in \mathbb{Q}, t \le x\}, \ \ x\in \mathbb{R}$.
\end{definition}

{\bf To prove:}
$b^r = \text{sup}(B(r)), \ \ r \in \mathbb{Q}$.

\begin{proof}
      Since $r$ is rational, $r-1$ is rational. $r-1 < r$, so $b^{r-1} \in B(r)$. Hence,
      $B(r)$ is nonempty.

      For all $b^t \in B(r)$ we have $t \le r$, which implies
      $b^t \le b^r$ (\ref{rationalExpComp}). Hence, $b^r$ is an upper bound of $B(r)$(which implies
      $\text{sup}(B(r))$ exists).

      Also, since $r \in \mathbb{R}$ and $r \le r$, $b^r \in B(r)$. Thus, no $b^t < b^r$
      can be an upper bound of $B(r)$.

      $\therefore \ \text{sup}(B(r)) = b^r$.
\end{proof}

\begin{definition}\label{realExpDef}
      $b^x \equiv \sup B(x), x\in \mathbb{R} $.
\end{definition}
\item

\begin{lemma}
      \label{ineqRoot}
      For all $x, y \in \mathbb{R}^+$ and $m, n \in \mathbb{Z}$, $m \ge 1$, $x^m < y^n$ implies
      $x < (y^n)^{1/m}$.
\end{lemma}

\begin{proof}
      It suffices to prove that $x \ge (y^n)^{1/m} \implies x^m \ge y^n$.
      \begin{align*}
                        & x \ge (y^n)^{1/m}                            \\
            \implies \  & x^m \ge ((y^n)^{1/m})^m & \ref{ineqMul}      \\
            \implies \  & x^m \ge y^n             & \text{Rudin, 1.21}
      \end{align*}
\end{proof}
\begin{lemma}
      \label{powerPartition}
      For all $b > 1$, the set $\{[b^n, b^{n+1})\ |\ n \in \mathbb{Z}\}$ is a partition
      of $\mathbb{R}^+$.
\end{lemma}
I have not been able to prove this rigourously.

\begin{lemma}
      \label{densityOfBr}
      For all $\alpha, \beta \in \mathbb{R}^+,\  \alpha < \beta$ there exists a number of the form
      $b^r$ where $b\in \mathbb{R},\  b > 1,\  r\in \mathbb{Q}$ such that
      $\alpha < b^r < \beta$.
\end{lemma}

\begin{proof}
      Since $\alpha < \beta,  \ \beta/\alpha > 1$. From \ref{powerPartition}, there
      exists $k \in \mathbb{Z}$ such that $b \in [(\beta/\alpha)^{k-1}, (\beta/\alpha)^{k})$.
      Thus, $b < (\beta/\alpha)^{k}$, or $\alpha^kb < \beta^k$. Since $b>1$, we have $k-1 \ge 0$
      or $k \ge 1$ (required for Lemma \ref{ineqRoot}).

      Also from \ref{powerPartition},
      there exists $l \in \mathbb{Z}$ such that $b^{l-1} \le \alpha^k < b^{l}$. Multiplying
      this inequality by $b$ yields $b^{l} \le \alpha^kb$.

      Chaining the above inequalities gives
      \begin{align*}
                     & b^{l-1} \le \alpha^k < b^{l} \le \alpha^kb< \beta^k &                \\
            \implies & \alpha^k < b^{l} < \beta^k                          &                \\
            \implies & \alpha < b^{l/k} < \beta.                            & \ref{ineqRoot}
      \end{align*}

\end{proof}

{\bf To prove:}
$b^{x+y} = b^xb^y \ \forall \,x, y \in \mathbb{R}$.

\begin{proof}
      Let $x+y \notin \mathbb{Q}$. This allows us to ignore the slackness of the inequality
      in Definition \ref{BxDef}. Let $x_1, x_2 \in \mathbb{Q}$ such that $x_1 < x < x_2$ (the existence of
      such $x_1$ and $x_2$ is guaranteed by the archimedean property). Similarly,
      let $y_1, y_2 \in \mathbb{Q}$ such that $y_1 < y < y_2$. Note that adding the two
      inequalities gives $x_1 + y_1< x + y< x_2 + y_2$.

      Now, since $x_1 < x$, it must be that $b^{x_1} \in B(x)$ by Definition \ref{BxDef}.
      So, $b^{x_1} \le \sup B(x)$, which implies $b^{x_1} \le b^x$ by Definition \ref{realExpDef}.
      Since $x_1 \ne x$, we have $b^{x_1} < b^x$. From \ref{rationalExpComp}, we have
      $b^{x_1}  < b^{x_2}$. It cannot be the case that $b^{x_2} < b^x$,
      since $x_2 \notin B(x)$. Thus, we have $b^{x_1}  < b^x<b^{x_2}$. Similarly,
      $b^{y_1}  < b^y<b^{y_2}$.
      \begin{align*}
                     & b^{x_1}  < b^x<b^{x_2} \text{ and } b^{y_1}  < b^y<b^{y_2}                 \\
            \implies & b^{x_1}  b^{y_1}  < b^xb^y<b^{x_2}b^{y_2}                  & \ref{ineqMul} \\
            \implies & b^{x_1+y_1}  < b^xb^y<b^{x_2+y_2}                          & \text{6(b)}
      \end{align*}

      This tells us that $b^xb^y$ is an upper bound for the set
      $S_1 \equiv \{b^t \ | \ t \in \mathbb{Q}, \ t < x + y\}$ and a lower bound for the set
      $S_2 \equiv \{b^t \ | \ t \in \mathbb{Q}, \ t > x + y\}$. Thus, it must be the case that
      \begin{align*}
            \sup{S_1} \le b^xb^y \le \inf{S_2}.
      \end{align*}
      $\sup{S_1}$, from Definition \ref{realExpDef}, is $b^{x+y}$. Thus, it suffices to
      prove that $\sup{S_1} = \inf{S_2}$. FTSOC, assume otherwise.
      \begin{enumerate}[(i)]
            \item
                  $\sup{S_1} < \inf{S_2}$\\
                  From \ref{densityOfBr}, there exists $b^r,\ r\in \mathbb{Q}$ such that
                  $\sup{S_1} < b^r < \inf{S_2}$. $b^r \notin S_1$, which implies that
                  $r \ge x+y$, which is equivalent to $r > x+y$ since $x+y$ is not rational.

                  Similarly, $b^r < \inf{S_2}$ implies that $b^r \notin S_2$, and $r \le x+y$
                  which is equivalent to $r < x+y$. \hfill  $\Rightarrow\Leftarrow$
            \item
                  $ \inf{S_2}< \sup{S_1}$\\
                  From the definition of supremum, there must exist $b^r \in S_1$
                  such that $ \inf{S_2}< b^r \le \sup{S_1}$. It follows that
                  $r < x+y$.

                  Similarly, from the definition of infimum, there must exist $b^{r'} \in S_2$
                  such that $\inf{S_2} \le b^{r'} < b^r$. It follows that $r' > x+y$.

                  $r' > x+y > r
                        \implies  r' > r$.

                  However, $b^{r'} < b^r \implies r' < r$.\hfill  $\Rightarrow\Leftarrow$
      \end{enumerate}
      Thus, $b^{x+y} = \sup{S_1} = \inf{S_2} = b^xb^y$.

      If $x+y \in \mathbb{Q}$, it suffices to prove that
      $\sup\{b^t \ | \ t \in \mathbb{Q}, \ t < x + y\} = \sup\{b^t \ | \ t \in \mathbb{Q}, \ t \le x + y\}
            = b^{x+y}$. This is trivially true from Lemma \ref{densityOfBr}
      ($\forall \,c < b^{x+y} \ \exists \,b^r$ such that $c < b^r <b^{x+y}$. $b^r$ is
      clearly a member of the first set. Hence $c$ is not an upper bound).
\end{proof}
\end {enumerate}


\section{}
 {\bf Definitions}\\
$b\in \mathbb{Q},\ b > 1$.\\
$y\in \mathbb{Q},\ y > 0$.

\begin{enumerate}[(a)]
      \item
            {\bf To prove:} $b^n-1 \ge n(b-1) \ \forall \, n \in \mathbb{Z}^+$.
            \begin{proof}
                  This can be proved using mathematical induction. For $n=1$,
                  $b-1 \ge b-1$. Let the inequality hold for some $k \in \mathbb{Z}^+$. Then,
                  \begin{align*}
                                 & b^k-1 \ge k(b-1)            \\
                        \implies & b(b^k -1) \ge k(b-1)        \\
                        \implies & b^{k+1} -1 \ge k(b-1) + b-1 \\
                        \implies & b^{k+1} -1 \ge (k+1)(b-1)
                  \end{align*}
            \end{proof}
      \item
            {\bf To prove:} $b-1 \ge n(b^{1/n}-1)\ \forall \, n \in \mathbb{Z}^+$.\\
            Plug $b^{1/n}$ in place of $b$ in 7(a).

      \item
            {\bf To prove:} $t>1\text{ and } n > (b-1)/(t-1) \implies b^{1/n} < t$.
            \begin{proof}
                  From (b), we have
                  \begin{align*}
                        \frac{b-1}n+1 \ge b^{1/n}.
                  \end{align*}
                  Rearranging the given inequality gives us
                  \begin{align*}
                        t > \frac{b-1}n+1.
                  \end{align*}
                  It follows that $t > b^{1/n}$.
            \end{proof}
      \item
            {\bf To prove:} $b^w < y \implies (\exists \,n \in \mathbb{Z}^+ \text{ such that } b^{w+1/n} < y$).
            \begin{proof}
                  Let $t = y/b^w$. $t>1$. Consider $n \in \mathbb{Z}^+$ such that $n > (b-1)/(t-1)$. The existence
                  of such an $n$ is guaranteed by the archimedean property. Then, by 7(c), we have
                  $b^{1/n} < y/b^w$, which gives us $b^{w + 1/n} < y$.
            \end{proof}
      \item
            {\bf To prove:} $b^w > y \implies (\exists \,n \in \mathbb{Z}^+ \text{ such that } b^{w-1/n} > y$).
            \begin{proof}
                  Let $t = b^w/y$. $t>1$. Consider $n \in \mathbb{Z}^+$ such that $n > (b-1)/(t-1)$. The existence
                  of such an $n$ is guaranteed by the archimedean property. Then, by 7(c), we have
                  $b^{1/n} < b^w/y$, which gives us $y < b^{w - 1/n}$.
            \end{proof}
      \item
            $A \equiv \{w\ |\ b^w<y\}$\\
            $x=\sup{A}$\\
            {\bf To prove:} $b^x = y$.
            \begin{proof}
                  FTSOC, assume otherwise.
                  \begin{enumerate}[(i)]
                        \item $b^x < y$\\
                              From 7(d), there exists $n \in \mathbb{Z}^+$ such that $b^{x+1/n} <y$,
                              which implies $x + 1/n \in A$. However, it is given that $x$ is an upper bound
                              of $A$. \hfill $\Rightarrow\Leftarrow$
                        \item $b^x > y$ \\
                              From 7(e), there exists $n \in \mathbb{Z}^+$ such that $b^{x-1/n} >y$.
                              Clearly, $x - 1/n$ is an upper bound of A.
                              However, it is given that $x$ is the smallest upper bound
                              of $A$. \hfill $\Rightarrow\Leftarrow$
                  \end{enumerate}
            \end{proof}
      \item
            {\bf To prove:} The $x$ obtained in (f) is unique.\\
            $x$ is defined as the supremum of $A$ in (f). The supremum of a set is always unique.
\end{enumerate}
\stepcounter{section}
\section{}
 {\bf Definitions}\\
$z = a + bi$\\
$w = c + di$\\
$<:\mathbb{C} \rightarrow \mathbb{C}: z<w \iff (a<c) \text{ or } (a = c \text{ and } b < d)$.

To show that this order turns $\mathbb{C}$ into an ordered set, we need to show trichotomy
and transitivity. The former is easily shown as follows.
\begin{enumerate}[(i)]
      \item $a < c \implies z < w$
      \item $a = c$ and
            \begin{enumerate}[(a)]
                  \item $d < b \implies z < w$
                  \item $d = b \implies z = w$
                  \item $d > b \implies w < z$
            \end{enumerate}
      \item $a > c\implies w < z$
\end{enumerate}

To show transitivity, consider $z, w, x(=e+fi) \in \mathbb{C}$, such that $z < w$ and $w < x$.
\begin{align*}
               & (a<c \text{ or } (a = c \text{ and } b < d))\text{ and } (c<e \text{ or } (c = e \text{ and } d < f)) \\ \\
      \implies & (a < c \text{ and } c < e) \text{ or }                                                                \\
               & (a<c \text{ and } (c = e \text{ and } d < f))\text{ or }                                              \\
               & ((a = c \text{ and } b < d) \text{ and } c<e)\text{ or }                                              \\
               & ((a = c \text{ and } b < d) \text{ and } (c = e \text{ and } d < f))                                  \\ \\
      \implies & (a < e) \text{ or }                                                                                   \\
               & (a<e \text{ and } d < f))\text{ or }                                                                  \\
               & (a<e \text{ and } b < d))\text{ or }                                                                  \\
               & (a = e \text{ and } b < f)                                                                            \\ \\
      \implies & (a < e) \text{ or }                                                                                   \\
               & (a = e \text{ and } b < f)                                                                            \\ \\
      \implies & z < x
\end{align*}
Complex numbers under lexicographic order have the smallest upper bound property.
consider the set $S = \{a+ib| a\in A, b\in B\}$ where $A$ and $B$ are sets of real numbers.
If $x = \sup{A}$ and $y = \sup{B}$, then $x+iy$ is clearly an upper bound of S. For any 
$z \in S, z<x+iy$, $z'$ can be found such that $z<z'<x+iy$ thanks to the smallest upper bound
property of the real numbers. 
\end{document}